
\documentclass[UTF8,a4paper,11pt]{ctexart}

% ============================================================
% 页面与基础宏包
% ============================================================
\usepackage[
    top=2.20cm,
    bottom=2.20cm,
    left=2.30cm,
    right=2.30cm,
    headheight=15pt
]{geometry}
\usepackage{amsmath}
\usepackage{booktabs}
\usepackage{longtable}
\usepackage{tabularx}
\usepackage{array}
\usepackage{enumitem}
\usepackage{fancyhdr}
\usepackage{lastpage}
\usepackage{hyperref}
\usepackage{xcolor}
\usepackage[most]{tcolorbox}
\hypersetup{hidelinks}

% ============================================================
% 黑白正式配色
% ============================================================
\definecolor{BoxGray}{HTML}{F5F5F5}
\definecolor{LineGray}{HTML}{B8B8B8}
\definecolor{TextGray}{HTML}{555555}

% ============================================================
% 正文格式
% ============================================================
\setlength{\parindent}{2em}
\setlength{\parskip}{0pt}
\linespread{1.18}

\setlist[itemize]{
    leftmargin=2em,
    itemsep=0.10em,
    topsep=0.20em,
    parsep=0pt
}
\setlist[enumerate]{
    leftmargin=2.2em,
    itemsep=0.10em,
    topsep=0.20em,
    parsep=0pt
}

\renewcommand{\arraystretch}{1.32}

% ============================================================
% 标题格式
% ============================================================
\ctexset{
    section = {
        format = \large\bfseries,
        beforeskip = 1.30em,
        afterskip = 0.55em
    },
    subsection = {
        format = \normalsize\bfseries,
        beforeskip = 0.95em,
        afterskip = 0.35em
    }
}

% ============================================================
% 页眉页脚
% ============================================================
\pagestyle{fancy}
\fancyhf{}
\fancyhead[L]{\small AI工具使用详情}
\fancyhead[R]{\small 2026年全国大学生数学建模竞赛}
\fancyfoot[C]{\small \thepage/\pageref{LastPage}}
\renewcommand{\headrulewidth}{0.4pt}
\renewcommand{\footrulewidth}{0pt}

% ============================================================
% 表格列
% ============================================================
\newcolumntype{Y}{>{\raggedright\arraybackslash}X}
\newcolumntype{L}[1]{>{\raggedright\arraybackslash}p{#1}}
\newcolumntype{C}[1]{>{\centering\arraybackslash}p{#1}}

% ============================================================
% 正式交互记录框
% ============================================================
\newtcolorbox{recordbox}[2]{
    enhanced,
    breakable,
    colback=white,
    colframe=LineGray,
    boxrule=0.65pt,
    arc=0pt,
    outer arc=0pt,
    left=4mm,
    right=4mm,
    top=3.2mm,
    bottom=3.2mm,
    before skip=0.7em,
    after skip=0.8em,
    title={
        \textbf{#1}
        \hfill
        \normalfont\small #2
    },
    coltitle=black,
    colbacktitle=BoxGray,
    fonttitle=\normalsize,
    boxed title style={
        boxrule=0pt,
        arc=0pt
    }
}

\newcommand{\recorditem}[1]{%
    \par\vspace{0.28em}
    \noindent\textbf{#1}\quad
}

% ============================================================
% 正文开始
% ============================================================
\begin{document}

% ============================================================
% 标题
% ============================================================
\begin{center}
    \vspace*{-0.7cm}

    {\zihao{2}\bfseries AI工具使用详情}

    \vspace{0.30cm}

    {\large 2026年全国大学生数学建模竞赛}

    \vspace{0.45cm}
\end{center}

\vspace{0.45cm}

本项目围绕含光伏、负荷、储能和外购电的微网调度问题展开。
建模过程中使用的生成式人工智能工具包括
\textbf{DeepSeek V4.1 Flash}、
\textbf{DeepSeek V4 Pro}
和
\textbf{GPT-image-2}。
其中，DeepSeek V4 Pro 主要用于需要较强综合推理的
备选方案比较、参数分析、信息可用范围检查和策略评价；
DeepSeek V4.1 Flash 主要用于代码实现、
局部逻辑排查以及论文文字审阅等高频辅助任务；
GPT-image-2 主要用于论文图示表达，
辅助生成论文总体路线图和问题二流程图的候选视觉方案。
团队负责数学模型建立、核心公式推导、参数最终确定、
实验运行、模型比较与最终选择、论文图示语义确定以及最终结论形成，
并对AI生成的代码承担最终核对与修改责任。
团队使用AI时，一般先由队员确定模型结构、程序流程或需要表达的图示逻辑，
再根据任务性质选择相应AI工具，
将具体问题、代码编写要求、公式、实验现象或流程图结构要求提交给AI辅助处理。
AI输出包括参考意见、代码片段初稿和候选视觉方案，
最终是否采用由团队根据题目条件、数学推导、程序运行、实验结果和论文实际结构判断。

% ============================================================
\section{AI工具使用概况}

本项目使用 DeepSeek V4.1 Flash、DeepSeek V4 Pro 和 GPT-image-2。
其中，需要综合比较、较长推理和建模判断的任务主要使用 V4 Pro；
代码片段与模块初稿的生成、代码检查、局部问题排查和文字审阅等任务主要使用 V4.1 Flash；
论文总体路线图和问题二流程图的视觉方案生成主要使用 GPT-image-2。

\begin{center}
\small
\begin{tabularx}{\textwidth}{
    L{0.22\textwidth}
    L{0.24\textwidth}
    Y
}
\toprule
\textbf{AI工具/模型} &
\textbf{主要使用阶段} &
\textbf{主要用途} \\
\midrule
DeepSeek V4 Pro &
建模方案比较、参数分析及策略评价 &
比较简单预测方法与复杂预测模型；
分析风险参数的统计含义；
检查预测数据和日内预报的可用范围；
讨论价格预测指标与最终调度费用之间的关系；
对关键建模选择提供备选分析思路。 \\
DeepSeek V4.1 Flash &
程序编写、调试、快速检查及论文整理 &
按团队给出的列结构与接口要求生成数据读取与结果写出的代码片段初稿；
检查 Python 程序中的时间索引、SOC递推、
滚动窗口和边界处理；
辅助排查局部代码逻辑；
审查论文中的结果表述、术语使用和因果措辞。 \\
GPT-image-2 &
论文图示设计与流程表达 &
根据团队已经确定的建模结构生成流程图候选方案；
辅助设计论文总体路线图和问题二流程图的布局、层级、箭头和信息分组；
为最终人工筛选、修订和定稿提供视觉参考。 \\
\bottomrule
\end{tabularx}
\end{center}

上述AI工具均未直接决定最终模型或最终方案。
对于涉及模型选择、参数确定和程序修改的关键内容，
团队均重新通过题目原文、数学推导和程序实验进行核验；
对于AI生成的代码片段，
团队逐段核对变量含义、下标与系数，
并用缩小的数值小例和全年数值门禁检查确认后才写入正式程序；
对于 GPT-image-2 生成的流程图候选，
团队逐项核对模块含义、箭头方向、因果关系和论文实际建模内容后再进行修改与定稿。
以下选取较有代表性的使用过程进行说明。

% ============================================================
\section{典型AI使用与交互记录}

% ============================================================
\begin{recordbox}
{交互记录 01：预测模型选择与复杂模型比较}
{问题二 · DeepSeek V4 Pro}

\recorditem{使用目的}
团队已经构造 F0、F1、F2 三种只使用目标日前历史数据的预测方法，
同时考虑是否进一步使用 LSTM、Transformer 等复杂时序模型。
需要判断复杂模型是否确有必要，
并避免因提前使用未来数据或不同预测方法的比较条件不一致而产生不公平比较。

\recorditem{主要提示内容}
团队向 DeepSeek V4 Pro 说明数据按日、按10分钟记录，
并明确预测时只能使用目标日之前已经获得的信息。
要求其比较简单预测方法和复杂时序模型时，
重点分析未来信息使用风险、过拟合风险、比较条件的一致性，
以及预测精度与最终调度费用之间的关系，
而不直接替团队确定最终模型。

\recorditem{AI反馈摘要}
DeepSeek V4 Pro 建议保证所有预测方法使用相同的历史信息，
并在相同数据区间内比较预测误差、应急购电量和最终购电费用。
同时指出，预测误差更低不一定意味着整个调度系统费用更低；
复杂模型只有在相同条件下能够稳定改善最终调度结果时，
其额外复杂度才具有实际价值。

\recorditem{团队处理}
团队采纳了“使用相同历史信息、在相同数据区间内比较，
并结合最终调度结果评价”的原则，
但没有直接采用AI对具体预测模型的选择意见。

\recorditem{人工核验}
团队重新逐日检查预测数据的截止时间，
并在统一历史信息和比较区间下运行不同预测方案。
比较预测误差、应急购电量和全年购电费用后发现，
复杂模型没有表现出稳定的最终费用优势，
因此最终没有采用 LSTM 或 Transformer 作为主预测模型。
\end{recordbox}

% ============================================================
\begin{recordbox}
{交互记录 02：残差分位参数与敏感性分析}
{问题二 · DeepSeek V4 Pro}

\recorditem{使用目的}
问题二利用历史净负荷预测误差的经验分位数确定风险裕量。
团队需要确定历史窗口 $W$ 和分位参数 $\alpha$，
并正确理解 $\alpha$ 的统计含义。

\recorditem{主要提示内容}
团队向 DeepSeek V4 Pro 询问
$\alpha=0.8$、$0.9$ 和 $0.95$
分别表示什么，
以及该参数是否能够直接理解为“整日可靠度”。
同时要求AI提出值得进一步检查的敏感性因素。

\recorditem{AI反馈摘要}
DeepSeek V4 Pro 指出，
$\alpha$ 对应历史预测残差分布中的经验分位位置，
不能直接解释为一天内全部时段同时满足要求的联合概率。
AI同时建议进一步检查储能效率、储能容量、
充放电功率、初始SOC、终端状态以及预测误差放大程度
对最终结果的影响。

\recorditem{团队处理}
团队采纳了AI对经验分位含义的解释，
并按照建议扩展敏感性分析；
但 $W$ 与 $\alpha$ 的最终取值没有直接采用AI建议，
而是由全年逐日模拟结果确定。

\recorditem{人工核验}
团队重新比较不同 $W$ 和 $\alpha$ 组合，
并分别改变储能效率、容量、功率、初始SOC和预测误差倍率。
最终依据全年购电费用、应急购电量及结果稳定性确定最终参数组合。
\end{recordbox}

% ============================================================
\begin{recordbox}
{交互记录 03：SOC递推与滚动优化程序排查}
{问题二至四 · DeepSeek V4.1 Flash}

\recorditem{使用目的}
滚动优化程序初版已经能够运行，
但团队检查代表日结果时发现，
部分时段的SOC变化与按递推公式推算的结果不一致，
因此需要进一步定位程序实现问题。

\recorditem{主要提示内容}
团队将SOC递推公式和相关 Python 代码片段提交给
DeepSeek V4.1 Flash。
当 $x_t$、$y_t$ 分别表示交流侧充电量和放电量时，
团队使用的SOC递推关系为
\[
S_{t+1}
=
S_t+0.9x_t-\frac{y_t}{0.9}.
\]
同时要求AI重点检查：
\begin{itemize}
    \item 充放电效率方向是否一致；
    \item 时间索引是否错位；
    \item 首时段执行是否正确；
    \item 滚动窗口末端是否存在边界问题；
    \item SOC达到上下限时如何修正充放电量；
    \item 应急购电是否可能被错误地用于主动充电。
\end{itemize}

\recorditem{AI反馈摘要}
DeepSeek V4.1 Flash 提示团队重点检查
滚动窗口边界数组与实际执行时段之间的对应关系，
并指出按照SOC上下限修正实际充放电量后，
可能出现“优化得到的充放电量”与“实际执行后的SOC变化”不一致的问题。

\recorditem{团队处理}
团队根据提示重新检查代码，
确认滚动边界数组和SOC边界处理确有需要修正的部分，
随后让该工具给出修正后的滚动边界处理与槽内裁剪代码，
团队按其自身的变量顺序和数据结构改写后接入程序，
并与修正前的结果逐槽比对。

\recorditem{人工核验}
修改后，团队把代表日的充放电计划导入电子表格，
按SOC递推公式逐槽独立复算储电量，
与程序输出逐槽比对，
随后逐时段检查母线能量平衡，
并对全年结果运行
\texttt{balance}、
\texttt{SOC chain}、
\texttt{overlap} 和
\texttt{settlement}
等程序检查。
复算结果与程序一致后才保留修改。
\end{recordbox}

% ============================================================
\begin{recordbox}
{交互记录 04：数据读取与结果写出代码生成}
{全文 · DeepSeek V4.1 Flash}

\recorditem{使用目的}
附件数据均为按日、
按10分钟排列的多张表格，
其中附件3的预报数据同时带有日期、
发布时刻和提前期三个维度；
正式结果还需要按官方模板逐行逐列回填到多张工作表中。
这类读入与写出代码工作量较大但逻辑固定，
团队希望把时间集中在模型与实验本身。
对应 \texttt{microgrid\_core.py} 的附件读入函数，
以及 \texttt{solution\_lib.py} 中的答案工作簿写出函数。

\recorditem{主要提示内容}
团队向 DeepSeek V4.1 Flash 说明各附件的列结构、
日期与时刻的对应关系以及模板工作表的行序与列序，
要求其生成读取附件数据、
把逐小时预报整理成按天与发布时刻存储的结构、
并按模板填写购电量、
充放电量和紧急购电量工作表的代码。

\recorditem{AI反馈摘要}
DeepSeek V4.1 Flash 给出按列切片读取数据、
按日期与发布时刻归档预报、
以及逐行回填模板的代码初稿，
并提示了空值处理与工作表名称匹配问题。
初稿默认附件3的日期列逐行完整填写，
与实际文件中日期列存在空白单元格的写法不完全一致。

\recorditem{团队处理}
团队保留按列切片读取与按模板回填的整体结构，
补充日期列向下填充、
形状与缺失值检查，
把读入封装为 \texttt{load\_inputs} 等函数，
把写出封装为 \texttt{\_fill\_plan\_sheet}、
\texttt{\_fill\_cycle\_sheet}、
\texttt{\_fill\_emergency\_sheet} 等函数，
供四个问题复用，
并统一各表列名与单位标注。

\recorditem{人工核验}
团队抽查若干天的原始表格与程序读入结果，
核对日期、
时刻与槽序号的对应关系；
随后对生成的答案工作簿逐表检查表名、
行列数和单元格内容是否与官方模板一致，
确认全年逐日数据无缺失后才用于正式结果输出。
\end{recordbox}

% ============================================================
\begin{recordbox}
{交互记录 05：日内预报更新与合同调整时序检查}
{问题三 · DeepSeek V4 Pro}

\recorditem{使用目的}
题目只在
0:00、6:00、12:00 和 18:00
提供新的光伏预测信息。
团队需要确认程序严格遵守四个规定的预报时点，
并避免已经执行的时段被重新优化或重复结算。

\recorditem{主要提示内容}
团队向 DeepSeek V4 Pro 说明四个预测发布时间和初版合同调整逻辑，
要求重点检查：
\begin{itemize}
    \item 程序是否在非规定预报时点读取新的预测信息；
    \item 是否存在使用未来信息的可能；
    \item 已经执行的时段是否保持不变；
    \item 当前实际SOC是否正确传递到下一次优化；
    \item 同一目标时段的合同差额是否可能重复结算。
\end{itemize}

\recorditem{AI反馈摘要}
DeepSeek V4 Pro 指出，
如果程序在每一个10分钟时段都根据最新信息重新生成剩余合同，
实际上相当于连续获得新的预测信息，
与题目只允许四次预报更新的条件不一致。
AI建议只在规定的预报时点更新购电计划：
每次仅在新的光伏预报发布后重新计算尚未执行时段的购电计划；
已经执行的时段不再调整，已发生费用不再改变；
剩余时段根据当前SOC和当前已知信息重新优化。

\recorditem{团队处理}
团队重新核对题目原文后确认这一问题存在，
因此否决连续更新的实现方式，
将程序改为仅在四个规定预报时点调整尚未执行时段的购电计划。

\recorditem{人工核验}
团队分别抽查
0:00、6:00、12:00 和 18:00
四个预报时点前后的数据版本、SOC和购电计划状态；
重新逐时段累计费用，
确认已执行时段不会被重新修改，
并检查每个目标时段只相对原合同进行一次最终差额结算。
此外，通过多组不同预报组合对比实验检查不同预报更新方式下的结果变化。
\end{recordbox}

% ============================================================
\begin{recordbox}
{交互记录 06：波动电价下的评价方式检查}
{问题四 · DeepSeek V4 Pro}

\recorditem{使用目的}
团队比较昨日同一时段、近7日均值、
近30日均值和近30日中位数等价格预测方式时发现，
价格 MAE、RMSE、尖峰误差与最终购电费用的排序并不完全一致，
因此需要确定合理的策略评价方式。

\recorditem{主要提示内容}
团队要求 DeepSeek V4 Pro 解释：
为什么价格预测误差更小，
不一定意味着实际购电费用更低；
同时要求其检查预测价格与实际结算电价在模型中的作用是否应严格区分。

\recorditem{AI反馈摘要}
DeepSeek V4 Pro 指出，
预测模型与决策模型的评价目标并不相同。
价格预测误差只反映预测值与实际电价之间的偏差，
而最终调度费用还受到误差发生时段、
储能状态、合同调整能力和应急购电等因素影响。
因此AI建议除 MAE、RMSE 外，
还同时比较尖峰价格误差、应急购电量和全年购电费用，
并将提前知道全天实际电价的情况仅作为不可实施的参考。

\recorditem{团队处理}
团队采纳多指标评价方式，
没有以价格MAE最小作为唯一的策略选择依据，
同时在程序中严格区分预测电价与实际结算电价。

\recorditem{人工核验}
团队重新检查价格数据时间标签及程序调用位置，
确认预测阶段没有读取未来实际电价；
随后对不同价格预测方法分别运行全年逐日模拟，
综合比较价格误差、尖峰误差、应急购电量和购电费用后确定最终方案。
\end{recordbox}

% ============================================================
\begin{recordbox}
{交互记录 07：结果表述与论文文字审查}
{全文 · DeepSeek V4.1 Flash}

\recorditem{使用目的}
部分实验结果之间存在明显相关关系，
但团队希望避免将特定实验条件下的现象写成过强的普遍因果结论。

\recorditem{主要提示内容}
团队将部分结果分析文字提交给 DeepSeek V4.1 Flash，
要求检查是否存在：
将相关关系直接写成因果关系、
将单组实验结果过度推广，
或者将预测指标与调度效果简单等同的情况。

\recorditem{AI反馈摘要}
DeepSeek V4.1 Flash 建议减少
“必然导致”“显著提高”等过强措辞，
将其改写为带有实验条件的描述，
例如“在当前策略和信息条件下费用下降”
或“该现象与某项变化同时出现”。

\recorditem{团队处理}
团队参考AI意见重新修改相关文字，
但没有直接复制AI生成的完整段落。

\recorditem{人工核验}
修改后的结论逐项与预报组合对比、
敏感性分析和实际实验结果对应，
确保论文表述没有超出实验结果可以支持的范围。
\end{recordbox}

% ============================================================
\begin{recordbox}
{交互记录 08：论文总体路线图与问题二流程图生成}
{全文与问题二 · GPT-image-2}

\recorditem{使用目的}
论文撰写过程中，团队需要将已经确定的整体建模路线以及问题二的
“预测--风险修正--优化调度--结果评价”过程转化为适合论文展示的流程图。
为了提高图示的层次性、可读性和版面表达效果，
团队使用 GPT-image-2 辅助生成两类流程图候选：
一张用于展示全文四个问题之间的总体研究路线，
另一张用于展示问题二从历史数据、预测、风险修正到滚动优化和全年评价的具体流程。

\recorditem{主要提示内容}
团队先根据论文实际模型整理流程图中必须出现的模块、
模块之间的箭头方向、信息传递关系以及不得添加的内容，
再将这些要求提交给 GPT-image-2。
总体路线图要求体现数据输入与预处理、确定性日前计划、
预测误差下的风险修正、日内预报驱动的合同调整、
波动电价扩展以及最终结果评价之间的递进关系。
问题二流程图要求重点体现：
\begin{itemize}
    \item 历史光伏、负荷等数据输入与预测；
    \item 净负荷预测误差统计与风险裕量构造；
    \item 含储能微网优化模型；
    \item 滚动执行、SOC状态更新与必要的应急购电；
    \item 全年逐日模拟及费用、应急购电量等指标评价。
\end{itemize}
同时要求保持论文已有的因果关系和箭头方向，
不得自行增加论文未采用的预测模型、优化机制或评价指标。

\recorditem{AI反馈摘要}
GPT-image-2 根据团队给出的文字结构生成了多种候选图示，
主要差异体现在横向或纵向布局、模块分组、箭头组织和信息密度等视觉表达方式。
这些候选结果帮助团队比较了不同版式下流程层级和关键路径的清晰程度，
但部分候选中仍存在文字过密、模块位置不理想
或自动补充非必要说明等问题。

\recorditem{团队处理}
团队没有直接将AI生成图片作为最终论文插图。
对于总体路线图，团队根据四个问题的实际递进关系重新调整模块顺序和连接方式；
对于问题二流程图，团队重点保留预测、风险修正、优化调度和结果反馈主链条，
删除自动生成但论文未实际使用的内容，
并修改文字、箭头和布局，使其与正文公式和程序实现保持一致。
最终两张图均经过人工筛选、重绘或修订后用于论文。

\recorditem{人工核验}
团队逐项检查两张流程图中的模块名称、输入输出、箭头方向、
因果关系和时序关系是否与论文正文一致，
并重点确认问题二图中预测信息、风险裕量、SOC更新、滚动优化
和结果评价之间不存在逻辑倒置或额外假设。
最终仅保留能够由论文模型、公式或程序流程直接支持的图示内容。
\end{recordbox}

% ============================================================
\section{AI建议的采纳、修改与否决情况}

为明确AI输出对最终方案的实际影响，
主要建议及团队处理情况汇总如下。

\begin{longtable}{
    L{0.31\textwidth}
    C{0.15\textwidth}
    L{0.43\textwidth}
}
\toprule
\textbf{AI建议或输出} &
\textbf{团队处理} &
\textbf{人工核验依据} \\
\midrule
\endfirsthead
\toprule
\textbf{AI建议或输出} &
\textbf{团队处理} &
\textbf{人工核验依据} \\
\midrule
\endhead

保证不同预测模型使用相同历史信息并在相同数据区间内比较 &
采纳 &
逐日检查数据截止时间，并在相同数据区间重新计算预测误差、
应急购电量和购电费用 \\
使用历史残差经验分位数进行风险修正 &
修改后采纳 &
重新比较不同 $W$ 和 $\alpha$，并进行多组敏感性分析 \\
使用 LSTM、Transformer 等复杂模型作为主预测器 &
验证后否决 &
在相同历史信息和全年逐日模拟条件下，
没有获得稳定的最终费用优势 \\
检查滚动优化中的SOC边界和时段边界处理 &
修改后采纳（含由AI给出的边界处理与裁剪代码） &
逐槽独立复算代表日储电量、
逐时段核对母线平衡，
并通过全年程序检查后确认修改有效 \\
由AI生成数据读取、结果表格与图件输出的代码初稿 &
修改后采纳 &
抽查原始附件与程序读入结果，
核对日期与槽序号对应关系，
并逐表检查答案工作簿的表名、行列数与全年数据完整性
（相应函数位于 \texttt{load\_inputs} 与 \texttt{solution\_lib.py}） \\
连续获取新光伏预报并持续调整合同 &
否决 &
与题目规定的四个预报更新时间不一致，
存在引入未来信息的风险 \\
仅根据价格MAE最小选择问题四策略 &
未采纳 &
价格预测误差与最终购电费用不完全等价，
需要结合应急购电量和最终购电费用评价 \\
使用过强的因果关系表述 &
修改 &
结合预报组合对比和敏感性分析重新核对后，
改为限定实验条件的结果表述 \\
GPT-image-2 生成的总体路线图和问题二流程图候选 &
修改后采纳 &
逐项对照论文正文、公式和程序流程，
核对模块含义、箭头方向、因果与时序关系，
删除论文未采用的自动补充内容后再定稿 \\
\bottomrule
\end{longtable}

\vspace{0.6em}

以上为本项目使用生成式人工智能工具的总体情况。
各项交互均以团队的模型推导、
程序检查与实验结果为准，
AI给出的方案、建议和程序片段在采用前都经过核对与验证；
模型的建立与修正、
参数的最终取值、
全部数值实验的运行以及论文结论的形成均由团队完成。

\end{document}
